\documentclass[12pt,a4paper]{article}

% ==================================================
% 1. CONFIGURATION
% ==================================================

% Enable / disable title page
\newif\ifshowtitlepage
\showtitlepagetrue
% \showtitlepagefalse

% Document information
\newcommand{\worktitle}{Sprint 1 Retrospective}
\newcommand{\worksubtitle}{NAO Robot Exercise System}
\newcommand{\teamname}{NAOMI}
\newcommand{\subject}{NAO Robot Exercise System}

\newcommand{\university}{Slovak University of Technology in Bratislava}
\newcommand{\faculty}{Faculty of Informatics and Information Technologies}

\newcommand{\academicyear}{2026/2027}
\newcommand{\location}{Bratislava}
\newcommand{\documentyear}{2026/2027}

% Team leaders
\newcommand{\leaderone}{Ing. Jakub Perdek, PhD.}
\newcommand{\leadertwo}{doc. Ing. Fedor Lehocki, PhD.}

% Team members
\newcommand{\memberone}{Bc. Akhmadinurov Amal}
\newcommand{\membertwo}{Bc. Danilov Denis}
\newcommand{\memberthree}{Bc. Dmitriev Anton}
\newcommand{\memberfour}{Bc. Gainetdinov Ruslan}
\newcommand{\memberfive}{Bc. Kuzin Timofei}
\newcommand{\membersix}{Bc. Tikhoblazhenko Alena}
\newcommand{\memberseven}{Bc. Zimnitski Aliaksei}

% Sprint and meeting information -- CHECK BEFORE MERGE
\newcommand{\sprintname}{Sprint 1}
\newcommand{\sprintperiod}{26 October -- 9 November 2026}
\newcommand{\meetingdate}{27 October 2026, 3 November 2026}
\newcommand{\meetingtime}{16:00}
\newcommand{\meetingplace}{NAO robot laboratory, FIIT STU}
\newcommand{\participants}{All team members: A. Akhmadinurov, D. Danilov, A. Dmitriev, R. Gainetdinov, T. Kuzin, A. Tikhoblazhenko, A. Zimnitski}
\newcommand{\reportauthor}{Bc. Gainetdinov Ruslan}


% ==================================================
% 2. PACKAGES
% ==================================================

\usepackage[
    a4paper,
    top=3cm,
    bottom=2.7cm,
    left=2.8cm,
    right=2.8cm,
    headheight=45pt,
    headsep=20pt
]{geometry}

\usepackage[T1]{fontenc}
\usepackage[utf8]{inputenc}
\usepackage[english]{babel}
\usepackage{lmodern}

\usepackage{graphicx}
\graphicspath{{../../assets/}{assets/}}
\usepackage[table]{xcolor}
\usepackage{fancyhdr}
\usepackage{titlesec}
\usepackage{setspace}
\usepackage{lastpage}
\usepackage{array}
\usepackage{float}
\usepackage{caption}
\usepackage{booktabs}
\usepackage{tabularx}
\usepackage{csquotes}

% No bibliography in progress reports, so biblatex is not loaded.

\usepackage[hidelinks]{hyperref}


% ==================================================
% 3. COLORS
% ==================================================

\definecolor{primary}{HTML}{000000}
\definecolor{secondary}{HTML}{000000}
\definecolor{muted}{HTML}{000000}
\definecolor{linecolor}{HTML}{000000}
\definecolor{naomiTableHeader}{HTML}{17445A}
\definecolor{naomiTableStripe}{HTML}{EDF7FA}
\definecolor{naomiTableRule}{HTML}{66BBD0}

\newcolumntype{L}[1]{>{\raggedright\arraybackslash}p{#1}}
\newcolumntype{Y}{>{\raggedright\arraybackslash}X}


% ==================================================
% 4. TYPOGRAPHY
% ==================================================

\setstretch{1.2}
\setlength{\parindent}{0pt}
\setlength{\parskip}{7pt}

\titleformat{\section}{\Large\bfseries\color{primary}}{\thesection}{0.8em}{}
\titleformat{\subsection}{\large\bfseries\color{primary}}{\thesubsection}{0.8em}{}
\titleformat{\subsubsection}{\normalsize\bfseries\color{primary}}{\thesubsubsection}{0.8em}{}

\titlespacing*{\section}{0pt}{20pt}{10pt}
\titlespacing*{\subsection}{0pt}{14pt}{7pt}
\titlespacing*{\subsubsection}{0pt}{12pt}{5pt}

\captionsetup{font=small, labelfont=bf, justification=centering}


% ==================================================
% 5. HEADER AND FOOTER
% ==================================================

\newcommand{\naomiheaderfooter}{%
    \fancyhf{}
    \fancyhead[L]{\includegraphics[width=4cm,height=1.15cm,keepaspectratio]{horizontal.png}}
    \fancyhead[R]{%
        \begin{tabular}{r}
            {\footnotesize\bfseries \worksubtitle}\\[3pt]
            {\scriptsize Academic Year \academicyear}
        \end{tabular}}
    \fancyfoot[L]{{\scriptsize \teamname\enspace|\enspace\subject}}
    \fancyfoot[R]{{\scriptsize \thepage\enspace/\enspace\pageref{LastPage}}}
    \renewcommand{\headrulewidth}{0.7pt}
    \renewcommand{\footrulewidth}{0.4pt}%
}

\pagestyle{fancy}
\naomiheaderfooter
\fancypagestyle{plain}{\naomiheaderfooter}


% ==================================================
% 6. TITLE PAGE
% ==================================================

\newcommand{\makemytitlepage}{
    \begin{titlepage}
        \thispagestyle{empty}
        \centering

        % University logo
        \vspace*{0.3cm}
        \includegraphics[width=6cm,height=5cm,keepaspectratio]{vertical.png}

        \vspace{0.5cm}
        {\large\bfseries \university\par}
        \vspace{0.2cm}
        {\normalsize \faculty\par}

        % Project and document title
        \vspace{1.7cm}
        {\small\MakeUppercase{\worksubtitle}\par}
        \vspace{0.5cm}
        \begin{minipage}{0.95\textwidth}
            \centering
            {\fontsize{25}{32}\selectfont\bfseries \worktitle\par}
        \end{minipage}

        % Sprint period and author
        \vspace{0.6cm}
        {\normalsize \sprintperiod\par}

        % Team information
        \vfill
        \begin{minipage}{0.9\textwidth}
            \centering
            \begin{minipage}[t]{0.46\linewidth}
                \centering
                {\normalsize\bfseries Team Leaders\par}
                \vspace{0.25cm}
                {\small
                \leaderone\par
                \leadertwo\par}
            \end{minipage}%
            \hfill
            \begin{minipage}[t]{0.48\linewidth}
                \centering
                {\normalsize\bfseries Team Members\par}
                \vspace{0.25cm}
                {\small
                \memberone\par
                \membertwo\par
                \memberthree\par
                \memberfour\par
                \memberfive\par
                \membersix\par
                \memberseven\par}
            \end{minipage}
        \end{minipage}

        % Bottom information
        \vfill
        {\small \location\enspace\textbullet\enspace\academicyear\par}
    \end{titlepage}
}


% ==================================================
% 7. IMAGES
% ==================================================

% \imageplaceholder{Text}{Caption} -- until the screenshot is added
\newcommand{\imageplaceholder}[2]{
    \begin{figure}[H]
        \centering
        \begingroup
        \setlength{\fboxsep}{0pt}
        \fbox{
            \begin{minipage}[c][5cm][c]{0.85\textwidth}
                \centering
                {\Large\bfseries IMAGE}
                \vspace{0.3cm}

                {\normalsize #1}
            \end{minipage}
        }
        \endgroup
        \caption{#2}
    \end{figure}
}

% \documentimage{path}{caption}{label}
\newcommand{\documentimage}[3]{
    \begin{figure}[H]
        \centering
        \includegraphics[width=0.85\textwidth,keepaspectratio]{#1}
        \caption{#2}
        \label{#3}
    \end{figure}
}


% ==================================================
% 8. TABLES
% ==================================================

\newenvironment{naomitable}[2]{%
    \begin{table}[H]
        \centering
        \caption{#1}
        \label{#2}
        \begingroup
        \small
        \renewcommand{\arraystretch}{1.4}
        \setlength{\tabcolsep}{6pt}
        \rowcolors{2}{naomiTableStripe}{white}
}{%
        \endgroup
    \end{table}
}


% ==================================================
% 9. DOCUMENT
% ==================================================

\begin{document}

\ifshowtitlepage
    \makemytitlepage
\fi


% ==================================================
% MEETING DETAILS
% ==================================================

\section{Meeting Details}

At the end of \sprintname{}, the \teamname{} team met to evaluate the sprint
in a Scrum retrospective. \sprintname{} was the first sprint in which the
team worked on the existing solution: it focused on the development
environment and on refactoring the project.

\begin{naomitable}{Retrospective meeting}{tab:meeting}
    \begin{tabularx}{\textwidth}{L{4.5cm} Y}
        \rowcolor{naomiTableHeader}
        \textcolor{white}{\bfseries Item} &
        \textcolor{white}{\bfseries Details} \\
        Date and time & \meetingdate, \meetingtime \\
        Place & \meetingplace \\
        Retrospective for & \sprintname{} (\sprintperiod) \\
        Participants & \participants \\
        \arrayrulecolor{naomiTableRule}\bottomrule[1pt]
    \end{tabularx}
\end{naomitable}


% ==================================================
% SPRINT OUTCOME
% ==================================================

\section{Sprint Outcome}

The team completed 12 of 14 tasks (36 of 44 story points). Two tasks were
not completed: \textbf{SCRUM-55 Set up CI pipeline to build and push Docker
images} and \textbf{SCRUM-58 Simplify and document environment
configuration}. Their implementation was finished, but the pull requests
were approved only on the day after the sprint ended, so both tasks were
moved to Sprint~2.

The main results of the sprint were working Docker images of all services,
smaller images, a clearer repository layout, the robot address and port
configurable through environment variables, a protected \texttt{master}
branch, automatic building and publishing of team documents, access to the
faculty server and the analysis of robot falls. The details are in the
\sprintname{} progress report.


% ==================================================
% MEETINGS IN THE SPRINT
% ==================================================

\section{Meetings in the Sprint}



\textbf{30 october 2026 -- supervisor meeting.}
The team presented the state of the existing solution. The team and the
supervisor agreed to refactor and optimize the project first and to start
the multi-state choreography epic only afterwards. Tasks SCRUM-58, 59 and 60
were added to the sprint after this decision.

\textbf{27 October 2026 -- team meeting.}
The team discussed why two pull requests from Sprint~0 were not reviewed
before the sprint ended: they were waiting over the weekend and nobody
checked them. The team agreed that open review requests must not be left
without attention at the end of a sprint. The rest of the meeting was used
to distribute the tasks of \sprintname{}.

\textbf{3 November 2026 -- team meeting.}
Each team member started the microservices and Choregraphe on their own
computer. Several members had problems connecting to the services or to the
robot. The team identified what had to be fixed so that the environment
works the same way for everyone.


% ==================================================
% COURSE OF THE MEETING
% ==================================================

\section{Course of the Meeting}

The meeting started with a short summary of \sprintname{}: the team went
through the completed tasks and the two tasks that were not completed.

\subsection{What Went Well}

\begin{itemize}
    \item The problems found when connecting to the robot in the Sprint~0
          retrospective (missing dependencies, hard-coded robot address)
          were fixed early in the sprint.
    \item Every change went through a pull request with a description and
          a review, as agreed in the Sprint~0 retrospective.
    \item Work on the services was split into small tasks, so several team
          members could work on the repository in parallel.
\end{itemize}

\subsection{What Was a Problem}

\begin{itemize}
    \item \textbf{Late estimation.} Tasks were estimated in story points only
          on 30 October, four days after the sprint started. Jira therefore
          shows a commitment of 0 and the burndown chart is not usable for
          the first days of the sprint.
    \item \textbf{Scope changes during the sprint.} Ten of the fourteen tasks
          were added after the sprint had started, most of them after the
          analysis of the existing solution and the supervisor meeting.
    \item \textbf{Slow reviews at the end of the sprint.} Two finished tasks
          missed the sprint because their pull requests were approved one
          day late.
    \item \textbf{Inconsistent methodologies.} The pull request approval rule
          is described differently in the Version Control Methodology, the
          Technical Methodology and the Definition of Done, and the use of
          the Testing status differs between the documents and practice.
    \item \textbf{Pull requests opened against the original repository.}
          Opening a pull request from the link printed by \texttt{git push}
          targets the supervisor's original repository instead of the team's
          fork; this happened several times.
\end{itemize}

% Optional: what each member said (what went well, what was a problem,
% what they would like to change).


% ==================================================
% IMPROVEMENTS
% ==================================================

\section{Changes for Sprint 2}

Based on the retrospective, the team introduced the following changes
from Sprint~2:

\begin{itemize}
    \item Tasks are estimated in story points during sprint planning,
          before the sprint starts.
    \item Pull requests are reviewed within two working days. In the last
          two days of the sprint, team members check open review requests
          every day.
    \item Tasks that arise during the sprint are added to the backlog first;
          they are added to the running sprint only if they block the sprint
          goal.
    \item The methodologies are aligned with the agreed rules (SCRUM-81 for
          pull request approvals, SCRUM-83 for the Testing status).
    \item Pull requests in \texttt{naoMicroservices} are opened through the
          compare page of the team's fork
          (\texttt{team26fiitstu/naoMicroservices}).
\end{itemize}

\end{document}